\documentclass{article}

% if you need to pass options to natbib, use, e.g.:
%     \PassOptionsToPackage{numbers, compress}{natbib}
% before loading neurips_2026

% The authors should use one of these tracks.
% Before accepting by the NeurIPS conference, select one of the options below.
% 0. "default" for submission
\usepackage{neurips_2026}

\usepackage[utf8]{inputenc}
\usepackage[T1]{fontenc}
\usepackage{hyperref}
\usepackage{bookmark}
\usepackage{url}
\usepackage{booktabs}
\usepackage{amsfonts}
\usepackage{amssymb}
\usepackage{nicefrac}
\usepackage{microtype}
\usepackage{xcolor}
\usepackage{amsmath}
\usepackage{amsthm}
\usepackage{graphicx}
\usepackage{array}
\usepackage{tabularx}
\makeatletter
\setlength{\@neuripsbelowcaptionskip}{10\p@}
\makeatother

\newcommand{\tv}{\mathrm{TV}}
\newtheorem{theorem}{Theorem}
\newtheorem{proposition}[theorem]{Proposition}
\newtheorem{corollary}[theorem]{Corollary}
\newtheorem{assumption}[theorem]{Assumption}

\title{ADSD: Prompt-Side Acceptance Collapse Attacks on Speculative Decoding}

\author{%
  Anonymous Authors \\
  Anonymous Institution \\
  Anonymous Address \\
  \texttt{anonymous@example.com} \\
}

\begin{document}

\maketitle

\begin{abstract}
Lossless acceleration schemes guarantee output equivalence, not operational
robustness. Speculative decoding is a particularly exposed example: its
throughput depends on fine-grained alignment between a small draft model and a
large target verifier, creating a prompt-controlled inference-cost surface. We
study adversarial prefixes that preserve user-visible task behavior while
collapsing verifier acceptance and erasing the latency benefit of speculative
decoding. We formulate this as a continuous-to-discrete prompt optimization
problem, show via a prompt-waywardness proposition that unconditional
projection guarantees are impossible, and prove under a TV-Lipschitz assumption
that the projected discrete attack gap is controlled by relaxed optimization
error and an explicit projection-distortion term. Algorithmically, we
instantiate this framework in a verifier-aligned straight-through optimizer
with soft-collapse, total variation, reverse KL, entropy, and
projection-distortion terms, and extend it to a calibration-batch
universal-prefix attack with multi-timestamp weighting and gradient-guided
discrete seed search. In a white-box setup with
\texttt{Qwen2.5-0.5B-Instruct-GPTQ-Int8} drafting for
\texttt{Qwen2.5-14B-Instruct-GPTQ-Int8}, Palmetto search reduces exact
projected first-token acceptance from $\alpha_1=0.1597$ for a manual seed to
$\alpha_1=0.0288$ for the optimized boundary attack. On a matched 263-sample
token-wise GSM8K benchmark, the boundary attack raises mean sample time from
$23.10$s to $29.81$s while reducing aggregate acceptance from $0.597$ to
$0.509$. The stronger universal-prefix attack raises mean sample time to
$38.61$s ($+67.1\%$ over the benign baseline), lowers aggregate acceptance to
$0.390$, lowers block efficiency to $3.731$, and preserves the benign GSM8K
accuracy exactly at $222/263$. These results expose a concrete vulnerability in
lossless acceleration: a short prompt can leave final answers intact while
forcing the verifier to do substantially more work.
\end{abstract}

\section{Introduction}

Speculative decoding accelerates large language model (LLM) inference by
letting a small draft model propose multiple tokens that are then verified by a
larger target model in parallel \citep{leviathan2023speculative}. The resulting
speedup is lossless at the output layer, but it is not free: it depends
critically on the degree of alignment between the draft distribution and the
target verifier. When draft-target agreement is high, long blocks of drafted
tokens are accepted and throughput improves substantially. When agreement
collapses, the verifier commits fewer drafted tokens, target-side correction
work increases, and the latency benefits of speculative decoding shrink.

Most prior work treats this acceptance bottleneck as an optimization problem to
be fixed. Recent systems improve the draft, adapt draft length, redesign
verification, or otherwise try to raise acceptance rates
\citep{liu2023online,goel2024direct,agrawal2024adaedl,bachmann2025judge,wang2025alignment,zhou2026hsd}.
This paper studies the complementary adversarial question: can an attacker
exploit the same acceptance bottleneck from the input side? In particular, can
a short prompt prefix preserve user-visible task behavior while forcing the
verifier into a substantially less efficient regime?

This attack surface is natural in deployment. The attacker does not need model
weights, scheduler privileges, or a custom decoding stack at inference time.
Instead, the attacker controls only the prompt text. Because speculative
decoding ties latency directly to the relation between the draft distribution
$q$ and the target distribution $p$, prompt-side inputs can act as
systems-level levers: they can induce overconfident draft proposals that the
target model repeatedly rejects, thereby converting speculative acceleration
into additional verifier work.

The technical obstacle is discrete. Gradient-based prompt optimization operates
most naturally over continuous token relaxations, but deployment requires a
discrete projected prefix. The main question is therefore not simply whether a
smooth surrogate can be written down, but when optimizing that surrogate yields
a strong \emph{discrete} attack after projection. This continuous-to-discrete
gap is the central methodological challenge for prompt-side attacks on
speculative decoding.

Our starting point is the token-wise verifier acceptance rule
\[
\alpha_i = \min\!\left(1, \frac{p(y_i \mid h_i)}{q(y_i \mid h_i)}\right),
\]
where $h_i$ is the speculative history and $y_i$ is the drafted token at step
$i$. Under draft sampling, the one-step expected acceptance satisfies
\[
\mathbb{E}[\alpha_i] = 1 - \tv(p_i, q_i),
\]
which immediately motivates an adversarial objective: maximize early
draft-target disagreement, since early verifier rejection is the most damaging
to end-to-end speculative efficiency. We formalize this intuition with a
continuous objective over relaxed prefixes, show that unconditional projection
guarantees are impossible via a prompt-waywardness proposition, and prove a
relaxation-gap bound under a TV-Lipschitz assumption that directly motivates a
projection-distortion penalty in the implemented optimizer.

On the algorithmic side, we instantiate this theory in two stages. The first
stage is a verifier-aligned soft-prefix optimizer that targets first-token
collapse at the assistant boundary using a straight-through relaxed prefix, a
soft-collapse term, TV, reverse KL, entropy sharpening, and an explicit
projection-distortion penalty. The second stage strengthens the attack into a
calibration-batch universal adversarial prefix optimized over target-generated
multi-timestamp rollouts. This universal-prefix pipeline combines a weighted
trajectory-level objective with a gradient-guided discrete seed-search stage,
allowing the optimizer to discover stronger discrete initializations before
continuous refinement.

We evaluate the attack in the reproduced HSD workspace using
\texttt{Qwen2.5-0.5B-Instruct-GPTQ-Int8} as the draft model and
\texttt{Qwen2.5-14B-Instruct-GPTQ-Int8} as the target model. The initial
Palmetto soft-prefix search consistently recovers the same five-token discrete
boundary attack across a sweep of projection-penalty strengths, reducing exact
first-token acceptance from $\alpha_1 = 0.1597$ for the manual comparison seed
to $\alpha_1 = 0.0288$ on the optimized search context. On a matched
263-sample token-wise GSM8K benchmark, that boundary attack raises mean sample
time from $23.10$s to $29.81$s while reducing aggregate acceptance from
$0.597$ to $0.509$ and block efficiency from $5.832$ to $4.961$.

The stronger result comes from the calibration-batch universal-prefix attack.
The best accuracy-preserving universal prefix raises mean sample time to
$38.61$s, lowers aggregate acceptance to $0.390$, lowers block efficiency to
$3.731$, and reduces decoding speed to $51.68$, while matching the benign
token-wise GSM8K accuracy exactly at $222/263$. In other words, a short
prompt-side prefix can preserve task accuracy while extracting a large
inference-cost penalty from a lossless speculative-decoding stack.

Our contributions are:
\begin{itemize}
\item We formulate prompt-side speculative-decoding slowdown as a
continuous-to-discrete adversarial optimization problem over verifier
acceptance rather than output semantics.
\item We prove a prompt-waywardness impossibility result and a relaxation-gap
bound that makes the projection distortion term in the optimizer theoretically
explicit.
\item We implement a verifier-aligned soft-prefix attack and a
calibration-batch universal-prefix extension with gradient-guided discrete
seed search, both of which are faithful to the released code path.
\item We show on a reproduced HSD/Qwen stack that prompt-only attacks can
materially degrade speculative efficiency, and that the strongest
accuracy-preserving universal prefix raises end-to-end runtime by $67.1\%$ on a
matched token-wise benchmark.
\end{itemize}

The broader message is that speculative decoding exposes a distinct
prompt-controlled systems vulnerability. Lossless output semantics do not imply
cost robustness. A user-controlled prefix can leave the final answer intact
while forcing the verifier to do substantially more work.

\section{Background and Preliminaries}
\label{sec:background}

Speculative decoding accelerates autoregressive generation by pairing a small
draft model with a large target verifier \citep{leviathan2023speculative}. Let
$q_i(\cdot \mid h_i)$ denote the draft distribution at speculative step $i$
given history $h_i$, and let $p_i(\cdot \mid h_i)$ denote the corresponding
target distribution. The draft proposes a token $y_i \sim q_i$, and the target
accepts it with probability
\[
\alpha_i = \min\!\left(1, \frac{p_i(y_i)}{q_i(y_i)}\right).
\]

The verifier either commits the drafted token or replaces it with a
target-sampled correction. Consequently, speculative efficiency depends
directly on the overlap between $p_i$ and $q_i$: high agreement yields long
accepted draft blocks, whereas low agreement forces frequent verifier
intervention.

This dependence can be written in a particularly convenient form. Under draft
sampling, the one-step expected acceptance satisfies
\[
\mathbb{E}[\alpha_i] = 1 - \tv(p_i, q_i),
\]
where $\tv$ is total variation distance. This identity makes the attack surface
explicit. Any prompt that increases early draft-target disagreement will, in
expectation, reduce verifier acceptance and therefore diminish speculative
speedup.

The systems metrics used throughout the paper are derived from the same
acceptance process. Let accepted drafted tokens denote the number of draft
tokens that survive verification, drafted tokens denote the total number of
tokens proposed by the draft, and target steps denote the number of verifier
steps taken by the target. Then
\[
\text{aggregate acceptance}
= \frac{\text{accepted drafted tokens}}{\text{drafted tokens}},
\]
\[
\text{block efficiency}
= \frac{\text{accepted drafted tokens}}{\text{target steps}}.
\]
Aggregate acceptance captures how much of the draft computation survives, while
block efficiency measures how much useful drafted work is recovered per
verifier step. These counter-based metrics complement wall-clock latency in the
experiments, and they are less sensitive to incidental runtime noise than raw
timing alone.

\section{Threat Model and Motivation}
\label{sec:threat}

We consider a white-box adversary that knows the target model, the draft
model, the tokenizer, and the deployed speculative verification rule. At
runtime, however, the adversary remains strictly prompt-bounded: the attacker
may modify only a short assistant-boundary prefix in the prompt. The attacker
cannot change model parameters, decoding temperature, token limits, system
prompts, verification logic, scheduler behavior, or hardware placement at
attack time.

The attacker seeks a discrete prefix $z \in \mathcal{V}^m$ that maximizes
speculative inefficiency while preserving user-visible task quality. At the
systems level, the objective is to increase the verifier work required per
request. We operationalize this through counter-based metrics such as aggregate
acceptance and block efficiency together with wall-clock latency. Formally, the
attacker seeks a short prefix that solves a constrained problem of the form
\[
\max_{z \in \mathcal{V}^m} \;\; \mathcal{L}_{\mathrm{cost}}(z)
\qquad
\text{s.t.}
\qquad
\mathcal{Q}(z) \geq \mathcal{Q}_{\mathrm{benign}} - \delta,
\]
where $\mathcal{L}_{\mathrm{cost}}$ denotes speculative inefficiency and
$\mathcal{Q}$ denotes user-visible task quality.

Why does such an attack matter? In multi-tenant LLM serving, speculative
decoding is deployed precisely because it increases throughput per GPU. When an
adversarial prefix collapses verifier acceptance, the system is forced back
toward a more sequential decoding regime. The immediate victim is the attacked
request, but the operational consequence is broader: extra verifier work
consumes shared accelerator time, reduces node throughput, and degrades the
quality of service seen by co-located users. In that sense, prompt-side
acceptance collapse functions as a resource-exhaustion attack on the
speculative serving stack, even though the adversary changes only the input
text and not the serving infrastructure itself.

The remaining challenge is methodological. The attacker needs a discrete prefix
that works in the deployed verifier, but the most powerful search tools are
continuous. The next section addresses that obstacle directly by formalizing
the projection gap and constructing an optimizer that treats it as part of the
attack objective rather than as an afterthought.

\section{Attack Methodology}
\label{sec:method}

\begin{figure}[h!]
\centering
\fbox{%
  \begin{minipage}[c][1.7in][c]{0.95\linewidth}
    \centering
    \textbf{Placeholder: ADSD Method Figure}\\[0.5ex]
    End-to-end attack pipeline. The final version should illustrate:
    (i) continuous prefix relaxation, (ii) projection-gap control with
    $\rho(U)$, (iii) verifier-aligned soft-prefix optimization, (iv)
    calibration-batch universal-prefix extension over target rollouts, and
    (v) discrete prefix evaluation in the speculative-decoding harness.
  \end{minipage}
}
\caption{Overview of the ADSD pipeline. Continuous prompt optimization is
coupled with explicit projection-gap control and verifier-aligned discrete
selection, then extended to a calibration-batch universal-prefix attack over
multiple prompts and timestamps.}
\label{fig:adsd-pipeline}
\end{figure}

\subsection{The Continuous-to-Discrete Projection Obstacle}

The attacker controls a discrete prefix, but gradient-based search is most
effective in a continuous relaxation of prompt space. Let
$E \in \mathbb{R}^{|\mathcal{V}| \times d}$ denote the shared token embedding
matrix. A relaxed prefix is
$U = (u_1, \dots, u_m)$, where each $u_t \in \Delta^{|\mathcal{V}|}$ is a
simplex vector over the vocabulary and the corresponding relaxed embedding is
\[
e_t(U) = E^\top u_t.
\]
This relaxation is attractive because it exposes gradient information over the
entire vocabulary at every position. The difficulty is that deployment still
requires a discrete projected prefix.

That obstacle is structural rather than incidental.

\begin{proposition}[Prompt Waywardness]
\label{prop:waywardness}
For any finite token embedding space with at least two tokens and any
$\delta > 0$, there exists a smooth objective $C$ over the convex hull of token
embeddings and a relaxed optimizer $U^*$ such that $U^*$ lies within distance
$\delta$ of a discrete token cell, yet the nearest-neighbor projection
$\Pi(U^*)$ is arbitrarily suboptimal among all discrete prefixes.
\end{proposition}

Prompt waywardness shows that proximity in embedding space does not imply
attack quality after projection. A successful continuous attack therefore
cannot rely on surrogate optimization alone; it must explicitly control the
relaxed-to-discrete gap.

\subsection{Bounding the Relaxation Gap}

Let $p_i^U$ and $q_i^U$ denote the target and draft next-token distributions at
speculative step $i$ under relaxed prefix $U$. We consider the weighted
disagreement objective
\[
C(U) = \sum_{i=1}^{\gamma} w_i \tv(p_i^U, q_i^U),
\]
where $w_i = \gamma - i + 1$ linearly emphasizes earlier speculative steps.
The corresponding discrete objective is
\[
C_d(z) := C(E(z)).
\]
To recover a deployable attack, define the token-wise nearest-neighbor
projection $\Pi(U) := (\Pi_1(U), \dots, \Pi_m(U))$ by
\[
\Pi_t(U) = \arg\min_{v \in \mathcal{V}} \|e_t(U) - E_v\|_2,
\]
and define the projection distortion
\[
\rho(U) := \sum_{t=1}^{m} \|e_t(U) - E_{\Pi_t(U)}\|_2.
\]

\begin{assumption}[TV-Lipschitzness]
\label{assump:tv-lipschitz}
For each speculative step $i$, there exist constants $L_i^p, L_i^q$ such that
for any two relaxed prefixes $U, V$,
\[
\tv(p_i^U, p_i^V) \le L_i^p \|U - V\|_{2,1},
\qquad
\tv(q_i^U, q_i^V) \le L_i^q \|U - V\|_{2,1},
\]
where
\[
\|U - V\|_{2,1} := \sum_{t=1}^{m} \|e_t(U) - e_t(V)\|_2.
\]
Let
\[
L_C := \sum_{i=1}^{\gamma} w_i (L_i^p + L_i^q)
\]
denote the aggregate Lipschitz constant of the weighted objective.
\end{assumption}

Under this assumption, the discrete quality of a projected relaxed solution can
be bounded explicitly.

\begin{theorem}[Relaxation Gap for Early Acceptance Collapse]
\label{thm:relax-gap}
Let $U^*$ be an $\varepsilon$-optimal solution of the relaxed problem such that
\[
C(U^*) \ge \sup_U C(U) - \varepsilon.
\]
Let
\[
z^* \in \arg\max_{z \in \mathcal{V}^m} C_d(z)
\]
be the optimal discrete prefix. Then
\[
C_d(\Pi(U^*)) \ge C_d(z^*) - \varepsilon - L_C \rho(U^*).
\]
\end{theorem}

\begin{proof}
By Assumption~\ref{assump:tv-lipschitz}, the objective $C$ is $L_C$-Lipschitz
under $\|\cdot\|_{2,1}$. Let $\hat{z} = \Pi(U^*)$. Since the continuous
embeddings of $\hat{z}$ are within $\rho(U^*)$ of $U^*$, we have
\[
C_d(\hat{z}) = C(E(\hat{z})) \ge C(U^*) - L_C \rho(U^*).
\]
Because the optimal discrete prefix is feasible in the relaxed space,
\[
C(U^*) \ge C_d(z^*) - \varepsilon.
\]
Combining the two inequalities yields the claim.
\end{proof}

We can also state an exact-recovery condition. Define the discrete attack
margin
\[
\Gamma := C_d(z^*) - \max_{z \neq z^*} C_d(z).
\]

\begin{corollary}[Exact Recovery under Discrete Margin]
\label{cor:exact-recovery}
Assume the optimal discrete adversarial prefix $z^*$ is unique
($\Gamma > 0$). If
\[
\varepsilon + L_C \rho(U^*) < \Gamma,
\]
then
\[
\Pi(U^*) = z^*.
\]
\end{corollary}

The theorem and corollary isolate the key design requirement: relaxed search
must improve the surrogate objective while also staying close to the discrete
token manifold. This is why ADSD uses an explicit projection-distortion
penalty rather than treating projection as a post-processing step.

\subsection{The ADSD Objective Formulation}

Motivated by Theorem~\ref{thm:relax-gap}, we construct a verifier-aligned
objective that couples disagreement maximization with explicit control of the
projection gap. The single-context attack targets the first verifier step,
since early rejection is the most damaging to token-wise speculative
throughput. Let $p^U$ and $q^U$ denote the target and draft next-token
distributions after the relaxed prefix $U$. ADSD optimizes
\begin{align*}
J(U) ={}&
\underbrace{\sum_{y} q^U(y)\,\mathrm{softplus}\!\left(\log q^U(y)-\log p^U(y)\right)}_{\text{soft collapse}}
+ \beta \,\tv(p^U, q^U)
+ \eta \,\mathrm{KL}(q^U \Vert p^U) \\
&- \tau \,\frac{1}{m}\sum_{t=1}^{m} H(u_t)
- \lambda \,\rho(U).
\end{align*}

The soft-collapse term rewards draft overconfidence on tokens that the target
assigns much smaller probability. The TV and reverse-KL terms enlarge the
surrounding disagreement region, the entropy term sharpens the relaxed simplex
vectors, and the projection penalty is the theory-guided component that keeps
the optimizer near a realizable discrete attack.

In the implementation, the prefix is parameterized by per-position vocabulary
logits. A softmax converts those logits into simplex vectors, and Adam updates
the logits directly. Each optimization step uses a straight-through argmax
estimator: the forward pass uses the hard projected token at each position,
while the backward pass differentiates through the softmax relaxation. The
released code therefore optimizes over the full vocabulary while keeping the
forward path faithful to the discrete token manifold used at deployment time.

Although the theory is phrased in terms of explicit nearest-neighbor
projection, the implementation uses argmax projection for computational
efficiency. The theoretical link is preserved by retaining the explicit
$L_2$ projection-distortion penalty $\rho(U)$ in the objective. Model
selection is verifier-aligned: although the optimizer follows the smooth
objective $J(U)$, the retained attack is ranked by its exact projected
first-token acceptance after hard projection in the actual
speculative-decoding pipeline.

In the experiments reported here, the core hyperparameters are a five-token
prefix, Adam learning rate $0.3$, temperature $1.0$, entropy weight $0.01$, TV
weight $1.0$, reverse-KL weight $0.1$, and gradient clipping at $1.0$. The
projection weight $\lambda$ is calibrated automatically from the ratio between
the unpenalized objective and the initial projection distortion, then swept on
Palmetto through a multiplicative factor. Because the search landscape is
highly non-convex, the implementation warm-starts the initial logits from a
short manual adversarial seed and searches the full vocabulary around that
initialization.

\subsection{Calibration-Batch Universal Prefix}

To move beyond a single boundary context, ADSD also optimizes a shared
universal prefix over a small calibration batch of GSM8K training prompts. For
each calibration prompt $x^{(j)}$, the target model first produces a fixed
continuation $y^{*(j)}_{1:K}$ of length $K$. These target-only rollouts define
the temporal contexts on which the universal prefix is optimized.

Let $p_i^{U,j}$ and $q_i^{U,j}$ denote the target and draft next-token
distributions at timestamp $i$ on calibration prompt $j$. The universal-prefix
objective is
\begin{align*}
J_{\mathrm{UAT}}(U) ={}&
\frac{1}{N}\sum_{j=1}^{N}
\Bigg[
\lambda_c \sum_{i=1}^{s} \bar{w}_i\,
\mathrm{SC}\!\left(q_i^{U,j}, p_i^{U,j}\right)
+ \beta \sum_{i=1}^{K} w_i\,\tv\!\left(q_i^{U,j}, p_i^{U,j}\right)
\\
&\qquad\qquad
+ \eta \sum_{i=1}^{K} w_i\,\mathrm{KL}\!\left(q_i^{U,j} \Vert p_i^{U,j}\right)
\Bigg]
- \tau \,\frac{1}{m}\sum_{t=1}^{m} H(u_t)
- \lambda \,\rho(U),
\end{align*}
where $w_i \propto K-i+1$ linearly emphasizes earlier timestamps and
$\bar{w}_i$ is the corresponding normalized restriction to the first $s$
timestamps. The term $\mathrm{SC}$ is the same soft-collapse signal used in the
single-context attack. The design preserves the theoretical emphasis on early
verifier failure while extending disagreement pressure across a full target
trajectory.

The universal-prefix implementation also includes a discrete initialization
stage before continuous refinement. Starting from the warm-start seed, ADSD
computes gradient-guided token replacements at each prefix position, retains
the top vocabulary candidates, and runs a small beam search over the prefix.
Candidate prefixes are ranked by exact projected first-token acceptance, with
weighted trajectory TV used as the tie-breaker. The best discrete candidate
then initializes the same straight-through optimization described above.

This hybrid discrete-plus-continuous pipeline is important in practice. It
places the continuous optimizer in a stronger basin, reduces dependence on a
single handwritten seed, and materially improves the quality of the projected
universal prefix recovered after refinement.

\section{Evaluation}
\label{sec:eval}

\subsection{Experimental Setup}

We evaluate ADSD in the reproduced HSD workspace of
\citet{zhou2026hsd} using \texttt{Qwen2.5-0.5B-Instruct-GPTQ-Int8} as the
draft model and \texttt{Qwen2.5-14B-Instruct-GPTQ-Int8} as the target model
under token-wise speculative verification. Unless otherwise noted, the
benchmark is GSM8K with the reproduced Qwen chat template and a matched
five-token benign prefix budget. All end-to-end numbers are reported against
the same benign token-wise baseline.

Each run uses batch size one, stochastic decoding
(\texttt{do\_sample=True}), and a per-sample cap of
\texttt{max\_new\_tokens=512}. We measure systems impact with wall-clock
latency, aggregate acceptance, block efficiency, and decoding speed. Latency
is recorded directly around the speculative \texttt{generate(...)} call.
Aggregate acceptance is the fraction of drafted tokens that survive
verification, and block efficiency is the number of accepted drafted tokens per
target verifier step. These counter-based metrics are computed from the saved
speculative traces in the reproduced HSD evaluator and therefore expose the
behavior of the verification mechanism independently of incidental timing
noise. Evaluating at batch size one follows standard speculative-decoding
practice \citep{leviathan2023speculative,zhou2026hsd}: the object of study is
the per-sequence verifier mechanism itself rather than downstream serving-stack
scheduling.

Before the end-to-end benchmarks, the Palmetto search consistently finds
projected prefixes that collapse early verifier acceptance. In the single-prompt
setting, the optimized five-token prefix reduces exact projected
$\alpha_1$ from $0.1597$ to $0.0288$. In the calibration-batch universal search
with $K=64$, every calibration size $N\in\{1,2,4,8\}$ yields a projected
prefix with calibration-batch $\alpha_1=0$, and the $N=4$ prefix delivers the
strongest accuracy-preserving end-to-end attack.

\subsection{End-to-End Latency Degradation}

ADSD substantially degrades verifier efficiency on unseen prompts while
preserving task accuracy. Table~\ref{tab:tokenwise-benchmark} reports the
matched 263-sample GSM8K benchmark comparing the benign boundary baseline with
the strongest accuracy-preserving universal prefix ($N=4$) and the strongest
cost-maximizing universal prefix ($N=2$).

\begin{table}[t]
\centering
\caption{Token-wise 263-sample GSM8K benchmark. Lower aggregate acceptance,
block efficiency, and decoding speed are worse for the system; higher mean
sample time is worse for the system.}
\label{tab:tokenwise-benchmark}
\footnotesize
\begin{tabular}{lccccc}
\toprule
Setting & Time (s) & Agg. acc. & Block eff. & Speed & GSM8K acc. \\
\midrule
Benign baseline & $23.10$ & $0.597$ & $5.832$ & $79.55$ & $222/263$ \\
Universal prefix ($N=4$) & $38.61$ & $0.390$ & $3.731$ & $51.68$ & $222/263$ \\
Universal prefix ($N=2$) & $51.60$ & $0.358$ & $3.506$ & $46.09$ & $202/263$ \\
\bottomrule
\end{tabular}
\end{table}

The systems signal is unambiguous. The $N=4$ universal prefix increases mean
sample time from $23.10$s to $38.61$s, a $67.1\%$ slowdown, while lowering
aggregate acceptance from $0.597$ to $0.390$ and block efficiency from $5.832$
to $3.731$. The same condition reduces HSD-style decoding speed from $79.55$
to $51.68$ while preserving GSM8K accuracy exactly at $222/263$. This
alignment between counter-based verifier statistics and wall-clock latency is
critical: it shows that the slowdown is a consequence of algorithmic verifier
degradation rather than an artifact of hardware noise. The $N=2$ universal
prefix pushes the cost attack further still, but at the price of task
degradation. For the core stealth claim of the paper, the $N=4$ prefix is the
most important result: a short prompt-side perturbation can deny the intended
speculative speedup while leaving user-visible task quality unchanged.

\subsection{Comparison with Discrete Search Baselines}

Bounded continuous optimization is strictly more effective than heuristic
discrete search for verifier manipulation. Table~\ref{tab:baseline-comparison}
compares ADSD against three discrete baselines: a pure random five-token
prefix, a GCG-adapted search baseline, and an AutoDAN-adapted search baseline.

\begin{table}[t]
\centering
\caption{Comparison with discrete search baselines on the token-wise GSM8K
benchmark. Search cost reports the dominant prompt-search workload required to
produce the final five-token prefix.}
\label{tab:baseline-comparison}
\footnotesize
\begin{tabular}{lccccc}
\toprule
Method & Search cost & Time (s) & Agg. acc. & Block eff. & GSM8K acc. \\
\midrule
Benign baseline & --- & $23.10$ & $0.597$ & $5.832$ & $222/263$ \\
Random prefix & $[XX]$ & $[XX]$ & $[XX]$ & $[XX]$ & $[XX]$ \\
GCG-adapted & $[XX]$ & $[XX]$ & $[XX]$ & $[XX]$ & $[XX]$ \\
AutoDAN-adapted & $[XX]$ & $[XX]$ & $[XX]$ & $[XX]$ & $[XX]$ \\
ADSD (UAT, $N=4$) & $[XX]$ & $38.61$ & $0.390$ & $3.731$ & $222/263$ \\
\bottomrule
\end{tabular}
\end{table}

The comparison isolates the key advantage of ADSD. Heuristic discrete methods
must search directly in token space, where verifier collapse is sparse and
projection failures are invisible until evaluation time. ADSD instead uses a
continuous verifier-aligned objective and a projection-aware regularizer to
reach deeper acceptance collapse with a short prompt budget. The resulting
prefixes therefore achieve stronger end-to-end degradation than discrete search
alone, while requiring materially less trial-and-error over the combinatorial
token space.

\subsection{Cross-Architecture Vulnerability}

The vulnerability is a property of speculative verification itself, not a quirk
of a single Qwen pairing. Table~\ref{tab:scaling-placeholder} evaluates ADSD on
larger Qwen and LLaMA draft-target pairs to test whether the same prompt-side
mechanism persists across architectures and scales.

\begin{table}[t]
\centering
\caption{Cross-architecture scaling of ADSD. Larger target models amplify the
absolute systems cost of verifier collapse because each rejected speculative
window forces more expensive target-side work.}
\label{tab:scaling-placeholder}
\footnotesize
\begin{tabular}{lcccc}
\toprule
Target $\rightarrow$ Draft & Condition & Time (s) & Block eff. & Accuracy \\
\midrule
Qwen2.5 14B $\rightarrow$ 0.5B & Benign & $23.10$ & $5.832$ & $222/263$ \\
Qwen2.5 14B $\rightarrow$ 0.5B & ADSD & $38.61$ & $3.731$ & $222/263$ \\
Qwen2.5 72B $\rightarrow$ 3B & Benign & $[XX]$ & $[XX]$ & $[XX]$ \\
Qwen2.5 72B $\rightarrow$ 3B & ADSD & $[XX]$ & $[XX]$ & $[XX]$ \\
LLaMA-3 8B $\rightarrow$ 1B & Benign & $[XX]$ & $[XX]$ & $[XX]$ \\
LLaMA-3 8B $\rightarrow$ 1B & ADSD & $[XX]$ & $[XX]$ & $[XX]$ \\
\bottomrule
\end{tabular}
\end{table}

The expected pattern is straightforward: if verifier collapse reflects a
fundamental weakness of draft-target acceleration rather than an idiosyncrasy of
one model family, then the same prompt-side mechanism should transfer across
architectures and become even more costly as target-model capacity grows. In
that regime, the absolute latency penalty becomes larger even when the relative
drop in acceptance is comparable.

\subsection{Ablation of Optimization Components}

The projection-aware design of ADSD is empirically necessary. Table
\ref{tab:ablation-placeholder} ablates the main terms of the optimization
objective, with special emphasis on the projection penalty $\rho(U)$$,$ which
is the direct algorithmic instantiation of the theoretical bound in
Section~\ref{sec:method}.

\begin{table}[t]
\centering
\caption{Ablation of ADSD optimization components. The full objective combines
early collapse pressure, trajectory-level distributional mismatch, and
projection-aware regularization.}
\label{tab:ablation-placeholder}
\footnotesize
\begin{tabular}{lcccc}
\toprule
Objective variant & Projected first token & Time (s) & Block eff. & Accuracy \\
\midrule
Full ADSD objective & $[XX]$ & $38.61$ & $3.731$ & $222/263$ \\
w/o projection penalty $\rho$ & $[XX]$ & $[XX]$ & $[XX]$ & $[XX]$ \\
w/o TV term & $[XX]$ & $[XX]$ & $[XX]$ & $[XX]$ \\
w/o reverse KL & $[XX]$ & $[XX]$ & $[XX]$ & $[XX]$ \\
w/o entropy term & $[XX]$ & $[XX]$ & $[XX]$ & $[XX]$ \\
\bottomrule
\end{tabular}
\end{table}

This ablation directly tests the central methodological claim of the paper. If
the relaxation-to-text gap were merely a nuisance detail, then removing the
projection-aware term would not materially affect the final discrete attack. Our
theory predicts the opposite: without a mechanism that keeps the relaxed prefix
close to the discrete manifold, optimization can find a strong continuous
adversarial state whose projected text no longer preserves verifier collapse.
The ablation therefore functions as an empirical validation of the
prompt-waywardness argument in Section~\ref{sec:method}.

\subsection{Discussion: Robustness Against Adaptive Defenses}

One natural defense is to detect persistent acceptance collapse and halt
drafting adaptively, as in designs such as AdaEDL. Such a fallback is useful,
but it does not eliminate the vulnerability studied here. The attacker does
not need to force the system into a permanently degraded regime to succeed; it
is already sufficient to deny the request the intended speculative speedup.
Moreover, adaptive fallbacks require an observation window before the system
can decide that drafting is no longer beneficial. The attack therefore still
extracts a penalty during the detection phase, after which the request is
reduced to ordinary autoregressive decoding rather than accelerated
speculative decoding. From the adversary's perspective, that outcome remains a
successful worst-case degradation of the algorithmic mechanism.


\section{Conclusion}

We introduced ADSD, a prompt-side attack view of speculative decoding centered
on the continuous-to-discrete relaxation gap. The theory section identifies the
projection obstacle, proves that unconditional nearest-neighbor recovery is
impossible, and shows that under TV-Lipschitz model behavior the projected
discrete attack is controlled by relaxed optimization error and projection
distortion. This directly justifies the manifold-attraction penalty in our
optimization objective. Empirically, the attack is not merely a first-token
curiosity: the optimized prefix family degrades end-to-end token-wise
speculative decoding on a matched 263-sample GSM8K evaluation. The initial
boundary attack already produces a substantial latency penalty, and the
calibration-batch universal-prefix extension amplifies that effect further:
runtime rises from $23.10$s to $38.61$s while aggregate acceptance and block
efficiency contract sharply, yet GSM8K accuracy remains at the benign baseline.
The main practical implication is clear. Lossless decoding does not imply cost
robustness. A prompt can preserve answer quality while forcing the verifier to
do substantially more work, turning speculative decoding into a
prompt-controlled inference-cost surface.

\appendix

\section{Implementation Notes}

Our local attack evaluation uses the same draft-target model pair as the HSD
reproduction workspace and the same Qwen chat template as the GSM8K evaluator.
The preliminary adversarial seeds were ranked by exact first-token acceptance
using white-box access to both models. A separate search utility in the
companion codebase evaluates prompt candidates and records the draft token,
target probability, and divergence statistics for each seed.


\bibliographystyle{plainnat}
\bibliography{references}

\input{checklist}

\end{document}
